problem:  
Let \( (M^n,g) \) be a compact, simply connected Riemannian manifold with **quasipositive curvature**, that is, \( \sec \ge 0 \) everywhere and \( \sec > 0 \) somewhere.  
Suppose \(M\) admits an **effective isometric torus action** \(T^r \curvearrowright M\) of rank  
\[
r \ge \left\lfloor \frac{n}{3} \right\rfloor.
\]
In the strictly positively curved setting, all known examples are **rationally elliptic**, and the **Bott–Grove–Halperin conjecture** predicts that every compact nonnegatively curved simply connected manifold is rationally elliptic.  

**Question:**  
Does the presence of **substantial symmetry** ensure rational ellipticity already under quasipositive curvature?  
Equivalently, can there exist a compact, simply connected Riemannian manifold with quasipositive curvature and an isometric torus action of rank at least \(\lfloor n/3\rfloor\) whose rational homotopy type is **hyperbolic** (i.e. with exponentially growing rational homotopy)?  

Why is it a "good" problem:  
- **Profound insight:**  
  This problem isolates a **symmetry‑restricted version** of the Bott–Grove–Halperin conjecture.  The conjecture asserts that rational ellipticity should hold for all compact nonnegatively curved spaces, but even the quasipositive case is entirely open.  By introducing a lower‑bound on symmetry rank—a setting in which geometric and topological constraints are most rigid—this question probes whether the conjectural ellipticity might be proved first in the high‑symmetry regime.  

- **Bridge between fields:**  
  To approach this problem one must combine tools from **Riemannian geometry** (curvature bounds and torus actions), **Alexandrov geometry** (study of orbit spaces and curvature of quotient metrics), and **rational homotopy theory** (ellipticity vs. hyperbolicity dichotomy and growth of homotopy groups). Establishing or refuting ellipticity under these assumptions would illuminate how curvature and symmetry interact to control algebraic topology.  

- **Extreme simplicity and conceptual depth:**  
  The statement is deceptively simple—"Does quasipositive curvature with sufficiently large symmetry force rational ellipticity?"—but it targets a major open conjectural mechanism linking curvature bounds with finiteness in rational homotopy theory.  Its resolution requires understanding the subtle propagation of partial curvature positivity through the orbit‑space geometry.  

- **Potential impact:**  
  A proof of ellipticity in this symmetry‑bounded quasipositive setting would provide the first significant verified instance of the Bott–Grove–Halperin conjecture beyond the known examples.  A counterexample, on the other hand, would demonstrate a genuine topological distinction between positive and quasipositive curvature.  In either case, the outcome would critically advance our understanding of the interplay between curvature, symmetry, and global topology.  

Thus, the problem achieves **conceptual clarity**, **links multiple major domains**, and **refines a central conjectural question** within a tractable geometric framework—making it a truly **good research problem**.